% ============================================================
% Section 12: Conclusion
% ============================================================

\section{Conclusion}
\label{sec:conclusion}

\subsection{Summary of Contributions}

We have presented \aletheion{}, a decoder-only LLM that integrates a
complete epistemic system as trainable neural modules. The main
contributions are:

\begin{enumerate}
    \item \textbf{Intrinsic epistemology}: To our knowledge, \aletheion{}
    is the first LLM to produce complete epistemic tomography (30+ fields)
    per token, trained end-to-end with the language model. While prior work
    has explored intrinsic uncertainty estimation \cite{gal2016dropout,
    lakshminarayanan2017simple}, these approaches are limited to scalar
    confidence or variance estimates and do not produce structured,
    multi-dimensional epistemic state per token.

    \item \textbf{5D Riemannian Manifold}: Tokens are mapped to a
    learned geometric space where geodesic distances define calibrated
    confidence.

    \item \textbf{Minimal overhead}: The 11 epistemic sub-heads add
    only {\raise.17ex\hbox{$\scriptstyle\sim$}}364K parameters ($<$0.3\%),
    demonstrating that epistemic capability does not require significant
    parameter increase.

    \item \textbf{Calibration by construction}: The MAD loss
    (\Cref{thm:calibration}) guarantees calibration at equilibrium,
    unlike post-hoc methods such as temperature scaling.

    \item \textbf{Collapse prevention}: The VI (\Cref{sec:vi}) and
    EidosDecay (\Cref{sec:eidos}) prevent manifold degeneration, a
    problem identified empirically in the predecessor system (ATIC).

    \item \textbf{Native Continual Learning}: EWC + Replay integrated
    into the trainer enable multi-phase training without catastrophic
    forgetting.

    \item \textbf{Scalability}: Configurations from 1M to 640B
    parameters with the same architecture, facilitating research
    (small models) and deployment (large models).
\end{enumerate}

\subsection{Comparison with ATIC}

\aletheion{} arose from the need to overcome the limitations of the
ATIC system, a deterministic orchestrator that wraps external LLMs:

\begin{table}[H]
\centering
\caption{ATIC vs \aletheion{} comparison.}
\label{tab:vs_atic}
\begin{tabular}{lcc}
\toprule
\textbf{Aspect} & \textbf{ATIC} & \textbf{\aletheion{}} \\
\midrule
Epistemic & Extrinsic & \textbf{Intrinsic} \\
Granularity & Per response & \textbf{Per token} \\
Calibration & Heuristic & \textbf{Via loss (BCE)} \\
Gradient & None & \textbf{End-to-end} \\
Dependency & External LLM & \textbf{Self-contained} \\
Latency & Multi-step pipeline & \textbf{1 forward pass} \\
Scalability & CPU-bound & \textbf{GPU, 1M--640B} \\
Continual Learning & None & \textbf{EWC + Replay} \\
\midrule
Governance & Complete & Pending \\
Agents & 9+ registered & Pending \\
Flexibility & Swappable base LLM & Fixed \\
\bottomrule
\end{tabular}
\end{table}

\subsection{Future Vision: Integration}

The ideal scenario is integrating \aletheion{} as the base LLM for ATIC,
combining:

\begin{itemize}
    \item Intrinsic epistemics (\aletheion{}) with orchestration (ATIC)
    \item Per-token tomography with governance and agents
    \item Calibrated confidence with policy binding
    \item Learned manifold with real-time MPC navigation
\end{itemize}

In this configuration, the 354M \aletheion{} would function as an
epistemic evaluator of larger LLMs' responses, adding calibrated
tomography to any external model.

\subsection{Future Work}

\begin{enumerate}
    \item \textbf{Empirical evaluation}: Calibration benchmarks (ECE,
    MCE, Brier score) against baselines (temperature scaling, ensemble,
    MC Dropout).

    \item \textbf{Transfer learning}: Verify whether the 5D manifold
    transfers across domains (e.g., code $\to$ scientific text).

    \item \textbf{Scale}: Train 1.3B and 7B versions to validate that
    the epistemic system improves with scale.

    \item \textbf{ATIC integration}: Replace ATIC's external LLM with
    \aletheion{}, evaluating impact on response quality and tomography
    precision.

    \item \textbf{Ablation study}: Individual impact of each loss
    component and sub-head on final calibration.

    \item \textbf{VI efficiency}: Compare end-to-end trained VI vs
    deterministic VI (ATIC) for collapse prevention.
\end{enumerate}

\subsection{Code and Reproducibility}

The complete source code, including all modules, tests (261), scaling
configurations, and training scripts, is available at:

\begin{center}
    \texttt{github.com/gnai-creator/aletheion-llm-v2}
\end{center}

The 354M parameter backbone is trained on 4$\times$ H100 in
{\raise.17ex\hbox{$\scriptstyle\sim$}}13 hours with 7B tokens from FineWeb-Edu,
followed by epistemic fine-tuning on a single RTX 4090 (24GB VRAM).
